\section{Résumé exécutif}
\label{sec:executive_summary}

\lettrine[lines=2, findent=2pt, nindent=0pt]{L}{a} \textbf{Stratégie FX Multi-Moteurs} est un portefeuille quantitatif qui combine quatre moteurs de rendement structurellement indépendants~: une composante de \emph{retour à la moyenne} intraday filtrée par un régime macroéconomique, une composante de \emph{suivi de tendance} journalier sur trois paires, une composante de \emph{retour à la moyenne} journalier sur l'indicateur RSI, et une composante de \emph{suivi de tendance} sur trois instruments directionnels (or, dollar-yen, argent). L'ensemble est piloté par une couche unifiée de ciblage de volatilité.

Le mandat retenu vise un rendement annualisé de l'ordre de \textbf{40\,\%}. Le plafond de drawdown qui figurait dans les versions antérieures de ce document a été \textbf{retiré à la demande du mandant}~: le risque n'est plus une contrainte du programme, il en est une conséquence mesurée. Ce document énonce donc ce que ce niveau d'ambition coûte en repli, sans l'atténuer.

\subsection*{Les chiffres qui font foi viennent du moteur qui exécute}

Deux moteurs mesurent ce portefeuille et ils ne mesurent pas la même chose. Le moteur de recherche (\texttt{vectorbtpro}) applique un levier au poids de position~; le moteur d'exécution (MetaTrader~5) ouvre un notionnel borné par la distance au stop. Les chiffres publiés ici sont ceux de MetaTrader~5, parce que c'est lui qui passe les ordres. Les métriques de recherche restent en annexe de robustesse, où elles servent à éprouver la structure de la stratégie, pas à annoncer une performance.

\begin{table}[H]
\centering
\caption{Résultats du portefeuille de production dans le testeur MetaTrader~5 sur EURUSD.c, 2021-01-01 à 2026-04-29. Ce sont les chiffres qui font foi~: ils proviennent du moteur qui exécute réellement les ordres, avec le \emph{spread} du courtier, l'arrondi des lots et les frais de portage.}
\label{tab:mt5_reference}
\begin{tabular}{lr}
\toprule
Métrique & Valeur \\
\midrule
Profit net & \$49,143 (dépôt \$10,000) \\
CAGR & \textbf{39.65\,\%} \\
Sharpe (MetaTrader~5) & 1.01 \\
Drawdown d'équité maximal & \textbf{47.04\,\%} \\
Facteur de profit & 2.18 \\
Nombre de transactions & 905 \\
\bottomrule
\end{tabular}
\end{table}


\begin{figure}[htbp]
\centering
\includegraphics[width=\linewidth]{mt5_balance_curve.png}
\caption{Solde du compte simulé dans le testeur MetaTrader~5, dépôt initial de 10\,000~USD, sur la totalité des données disponibles. La courbe intègre le \emph{spread} du courtier, l'arrondi des lots et les frais de portage.}
\end{figure}

\begin{insightbox}
\textbf{Le résultat est concentré, et c'est le fait le plus important du document.} La sleeve momentum --- or, dollar-yen et argent, équipondérés --- représente 15\,\% de l'allocation et 87,0\,\% du résultat net, pour 89~transactions sur 905. À l'intérieur de cette sleeve, \textbf{une position unique sur l'argent apporte 46,6\,\% du résultat net du portefeuille}, et les trois plus grosses en apportent 80\,\%. La performance annoncée dépend donc d'un très petit nombre de mouvements, sur une période exceptionnellement favorable aux tendances des métaux et du yen. Aucune moyenne, aucun ratio de Sharpe ne rend visible cette dépendance --- la section~\ref{sec:limitations} en tire les conséquences.

\textbf{Le poids de cette sleeve a été réduit le 28~juillet~2026}, de 20\,\% à 15\,\%~: le test de sur-ajustement appliqué à la grille d'allocations échoue (section~\ref{sec:advanced_robustness}), ce qui indique que le poids le mieux classé sur une moitié des données ne se transporte pas sur l'autre. La conséquence était écrite avant le calcul.
\end{insightbox}

\begin{table}[H]
\centering
\caption{Décomposition du résultat par sleeve. « Gold Momentum » en concentre 87.0\,\% pour 89~transactions~: le résultat du portefeuille dépend d'un petit nombre de positions, ce qu'aucune moyenne ne montre.}
\label{tab:mt5_by_sleeve}
\begin{tabular}{lrrrr}
\toprule
Sleeve & Transactions & Résultat net & Part & Taux de gain \\
\midrule
Gold Momentum & 89 & \$42,755 & 87.0\,\% & 41.6\,\% \\
TS Momentum & 479 & \$3,613 & 7.4\,\% & 62.8\,\% \\
MR Macro & 304 & \$1,515 & 3.1\,\% & 57.9\,\% \\
RSI Daily & 33 & \$1,261 & 2.6\,\% & 66.7\,\% \\
\bottomrule
\end{tabular}
\end{table}


\subsection*{Quatre conclusions pour la décision}

\begin{enumerate}
\item \textbf{Le mandat de rendement est tenu} sur la période mesurée~: 39,7\,\% de croissance annualisée, pour un profit net de 49\,143~USD sur un dépôt de 10\,000~USD.
\item \textbf{Le repli maximal d'équité atteint 47,0\,\%.} C'est le prix direct du niveau de risque choisi~: le dimensionnement des positions a été calibré pour atteindre la cible de rendement, et il multiplie le risque dans la même proportion.
\item \textbf{La performance repose sur une sleeve.} Retirer la sleeve momentum ramènerait le résultat à moins d'un dixième de son niveau. La diversification annoncée par l'architecture à quatre moteurs ne s'est pas traduite par une diversification du résultat.
\item \textbf{Une année sur six est négative} (2023, $-8{,}1$\,\%), et le repli d'équité dépasse celui du solde sur toute la période --- signe que des positions restent longtemps ouvertes en territoire défavorable.
\end{enumerate}

\begin{table}[H]
\centering
\caption{Résultat par année civile, mesuré sur le solde du compte. Le repli indiqué est celui du \emph{solde}~: il ne descend pas aussi bas que le repli d'équité, qui compte aussi les positions encore ouvertes.}
\label{tab:mt5_yearly}
\begin{tabular}{lrrrr}
\toprule
Année & Solde de fin & Rendement & Repli maximal & Transactions \\
\midrule
2021 & \$10,408 & 4.1\,\% & -5.3\,\% & 97 \\
2022 & \$16,759 & 61.0\,\% & -3.7\,\% & 199 \\
2023 & \$15,402 & -8.1\,\% & -15.6\,\% & 203 \\
2024 & \$19,827 & 28.7\,\% & -13.0\,\% & 181 \\
2025 & \$27,529 & 38.8\,\% & -7.4\,\% & 174 \\
2026 & \$59,143 & 114.8\,\% & -13.3\,\% & 51 \\
\bottomrule
\end{tabular}
\end{table}


Les sections suivantes développent le mandat, l'architecture du portefeuille, les quatre moteurs de rendement, la couche de risque, les résultats de recherche avec leurs tests de robustesse, l'écart entre recherche et exécution, et les conditions de déploiement.
